\section{Transport of the degree-three primitive}\label{sec:degree-three}

The projective presentation of $\mathcal M_n$ isolates a
one-dimensional Galois-invariant subspace of its degree-three
cohomology. The Drinfeld presentation shows that the linear trace
class has nonzero image in that subspace. This identifies the two
primitive classes up to a nonzero scalar; the rank-one filtration
then transports the identification from rank three to rank two.
All cohomology uses the integral complex followed by inversion of
$p$, as defined in \cref{sec:cohomology}.

\subsection{The integral projective bundle splitting}

\begin{lemma}\label{lem:projective-bundle}
Let $X$ be a small v-stack over $\Spd C$, and let $\mathcal W$ be a
rank-$(r+1)$ vector bundle on its untilt, with linear v-descent.
For $g:\mathbb P(\mathcal W)\to X$, the compatible classes
$h_m=c_1(\cO(1))\in H_v^2(\mathbb P(\mathcal W),(\Z/p^m)(1))$
give equivalences
\begin{equation}\label{eq:projective-finite}
 \lambda_m:
 \bigoplus_{j=0}^r(\Z/p^m)(-j)[-2j]
       \xrightarrow{\ \sim\ } Rg_*(\Z/p^m).
\end{equation}
Consequently
\begin{equation}\label{eq:projective-integral}
 \bigoplus_{j=0}^r
  \Rlim_m R\Gamma_v(X,(\Z/p^m)(-j))[-2j]
 \xrightarrow{\ \sim\ }
  \Rlim_m R\Gamma_v(\mathbb P(\mathcal W),\Z/p^m).
\end{equation}
The term $j=0$ is the original pullback.  Thus pullback is split
injective on $H^*_{\mathrm{int}}$ and on $H^*_{\mathrm b}$.
These statements are equivariant for actions on the vector bundle,
including arithmetic Galois actions.
\end{lemma}

\begin{proof}
Define $\lambda_m$ by the powers $h_m^j$, using adjunction and cup
product.  It is therefore a global map of derived pushforward objects,
compatible with coefficient reduction and all descent data.
For arithmetic descent, the linear action on $\mathcal W$ induces
descent on the actual line bundle $\cO(1)$.  Naturality of the Kummer
boundary makes $h_m$ invariant in its twisted coefficient group.

Over $\operatorname{Spa}(C,\cO_C)$, the finite-constant
projective-space calculation has basis
$1,h_1,\ldots,h_1^r$: apply
\cite[Construction~3.2.1 and Proposition~3.2.2]{LRZ2024}
with $\F_p$ coefficients, since $p$ is invertible on this base.
The first Chern class there is the Kummer boundary
\cite[Variant~3.1.8]{LRZ2024}, so this basis uses our actual classes.
To extend the calculation to arbitrary perfectoid test bases, apply
the integral Riemann--Hilbert functor
$\operatorname{RH}_{\Z_p}$ of \cite[Theorem~5.1.7]{ALM2026}.
Smooth proper primitive comparison
\cite[Proposition~5.2.2 and Example~5.2.3]{ALM2026} applies to
$\mathbb P_C^r\to\Spd C$ and the dualizable coefficient object
$\F_p$. Its mod-$p$ target is
\[
 \mathcal D_{\mathrm{FF}}(-,\F_p)
   =\mathcal D_{\widehat{\square}}(\cO^\flat)^{\varphi},
\]
by \cite[Remark~5.1.9 and the beginning of \S5.2]{ALM2026}.
The symmetric monoidal and pullback compatibility of the functor
carries the Kummer basis to its coefficient images. It therefore
gives the Frobenius-equivariant Chern-class equivalence with
$\cO^\flat$ coefficients over $C$.

After trivializing $\mathcal W$ on a perfectoid test object, its
projective bundle is a base change of $\mathbb P_C^r\to\Spd C$.
This projective morphism is $p$-bounded. Base change for solid
pushforward, \cite[Proposition~4.2.7(ii)]{ALM2026} with
$\Lambda=\F_p$, transports the equivalence to that test object.
The globally defined Chern-class map makes these local equivalences
compatible on overlaps, so they descend.

We identify the resulting solid calculation with ordinary v-sheaf
pushforward by testing on a perfectoid $T\to X$. Apply tensor-unit
Hom in the solid formalism on $T$. Adjunction and the tensor-unit
cohomology identification
\cite[Theorem~3.2.1(iii)--(iv)]{BMR2026} give
\[
 \bigoplus_{j=0}^r
    R\Gamma_v(T,\cO^\flat(-j))[-2j]
 \xrightarrow{\ \sim\ }
    R\Gamma_v(\mathbb P(\mathcal W)_T,\cO^\flat).
\]
This is natural in $T$ and induced by the same Chern classes.
It thus identifies their coefficient-image map with the equivalence
of ordinary derived v-sheaf pushforwards
\[
 \bigoplus_{j=0}^r\cO^\flat_X(-j)[-2j]
       \xrightarrow{\ \sim\ } Rg_*\cO^\flat.
\]

The displayed map is induced by the images of the actual mod-$p$
\'etale classes $h_1^j$.  It therefore commutes with Frobenius on
$\cO^\flat$, including its descent data.  Take the fiber of
$F-1$ on both sides.  The Artin--Schreier sequence on the v-site
identifies this fiber with $\F_p$; with twists it identifies the
$j$th fiber with $\F_p(-j)[-2j]$.
Since derived pushforward preserves fibers, the resulting equivalence
is exactly $\lambda_1$.
This step uses the Frobenius compatibility of the Chern-class map.

The Kummer classes $h_m$ are compatible under reduction.
The coefficient sequences
\[
 0\longrightarrow\Z/p^{m-1}\xrightarrow{\,p\,}\Z/p^m
       \longrightarrow\F_p\longrightarrow0
\]
give compatible exact triangles on the source and target of
$\lambda_m$.  Cup product with $h_m^j$ commutes with both coefficient
maps.  Induction on $m$ therefore proves
\cref{eq:projective-finite}.
Apply global sections and then the derived inverse limit.
The finite direct sum commutes with both operations and gives
\cref{eq:projective-integral}.  Projection to its zeroth summand
is a left inverse to pullback.  Exact inversion of $p$ preserves that
left inverse.
\end{proof}

\subsection{Galois weights identify the primitives}

Recall from \cref{prop:group-classes} the nonzero classes
\[
 e_{P_n}\in H^3_{\mathrm b}(P_n),\qquad
 e_{J_n}\in H^3_{\mathrm b}(J_n),\qquad n=2,3.
\]
For a class $e\in H^i_{\mathrm b}(A)$ and a classifying map
$f:Z\to B_CA$, write
\[
 e|_Z=f^*u_{B_C,A,\mathrm b}(e).
\]
Thus the class is first mapped from condensed group cohomology to the
relative geometric classifying stack by
\eqref{eq:finite-constant-cochains}, and then pulled back.  For a group
homomorphism $\rho:D\to A$, naturality gives
\[
 d^*u_{B_C,D,\mathrm b}(\rho^*e)
       =(B_C\rho\circ d)^*u_{B_C,A,\mathrm b}(e).
\]
We abbreviate the left side as $d^*\rho^*e$.  The associated framing
torsors on $Y_C$ give exactly the convention for $e_{P_n}|_{Y_C}$
fixed in \cref{sec:completed-tower}, by the same finite-level naturality.

\begin{proposition}\label{prop:primitive-comparison}
For $n=2,3$, the pullbacks of $e_{P_n}$ and $e_{J_n}$ to
$H^3_{\mathrm b}(\mathcal M_n)$ are nonzero.
There is a scalar $\kappa_n\in\Q_p^\times$ such that
\begin{equation}\label{eq:primitive-comparison}
 e_{P_n}|_{\mathcal M_n}
        =\kappa_n\,e_{J_n}|_{\mathcal M_n}.
\end{equation}
\end{proposition}

\begin{proof}
Use the common arithmetic action $G_F$ from
\cref{prop:two-realizations}.
The projective presentation is a projective bundle over $B_CP_n$.
Apply \cref{lem:projective-bundle} first to this relative geometric
base.  Its two degree-three summands use the finite coefficients
$\Lambda_m$ and $\Lambda_m(-1)$ on $B_CP_n$, respectively, before
their derived limits; higher projective summands would have negative
cohomological degree and vanish.  The compact geometric-point comparison
\eqref{eq:geometric-point-constants}, applied in the $P_n$ bar degrees
and with its Tate fiber, identifies these summands with the group
cohomology groups in the $G_F$-equivariant direct sum
\begin{equation}\label{eq:projective-degree-three}
 H^3_{\mathrm b}(\mathcal M_n)
 \cong H^3_{\mathrm b}(P_n)
          \oplus H^1_{\mathrm b}(P_n)(-1).
\end{equation}
The first summand is the image of the actual constant-cochain pullback,
hence contains a nonzero $e_{P_n}$.  The comparison is natural for the
Galois action on the $\Spd C$ base.  It identifies the untwisted group
cochain factors with trivial Galois action because Galois fixes $P_n$
and the untwisted constants.
On the second summand it acts by $\chi_{\mathrm{cyc}}^{-1}$.
As $\chi_{\mathrm{cyc}}(G_F)$ is open, this summand has no invariants.
Thus
\begin{equation}\label{eq:invariant-line}
 H^3_{\mathrm b}(\mathcal M_n)^{G_F}
       =\Q_p\,e_{P_n}|_{\mathcal M_n}.
\end{equation}

It remains to prove that the pullback of $e_{J_n}$ is nonzero.
Choose a compact open uniform subgroup $U\subset J_n$.
Restriction sends $e_{J_n}$ to a nonzero element of
$H^3_{\mathrm b}(U)$ by \cref{prop:group-classes}.
We will detect it after the further pullback to $[\Omega_C^{n-1}/U]$.

Let $K_n=\mathscr C_{\Omega_C^{n-1}}$ be the derived solid integral
cohomology complex formed with profinite parameters.
Descent, as in \cref{lem:solid-resolutions}, gives the first-quadrant
spectral sequence
\begin{equation}\label{eq:drinfeld-descent}
 E_2^{s,j}=H^s(U,H^j(K_n))(*)
 \ \Longrightarrow\
 H_{\mathrm{int}}^{s+j}([\Omega_C^{n-1}/U]).
\end{equation}
The complex $K_n$ is bounded below, and the group degree has the finite
bound supplied by the completed resolution.  The sequence therefore
has a finite filtration in every degree and may be rationalized
exactly.  The additional geometric row bound uses only point values:
\cref{thm:drinfeld} gives $H^j(K_n)(*)=0$ for $j>n-1$;
the finite Hom calculation below then makes every ordinary
point-valued $E_2^{s,j}$ in that row zero.  No vanishing assertion for
the entire condensed module $H^j(K_n)$ is needed.

Choose the finite projective resolution
$P_\bullet\to\Z_p$ over $\Lambda=\Z_p[[U]]$ from
\cref{lem:solid-resolutions}.  The unit
$\underline{\Z_p}\to K_n$ factors through $H^0(K_n)[0]$, since
$K_n$ is concentrated in cohomological degrees at least zero.
The degree-zero calculation and local constants argument following
\cref{thm:drinfeld} show that this unit is an isomorphism on point
values.  Each $\Hom_\Lambda(P_i,-)$ is a retract of a finite sum of
point values, and the equivariant unit commutes with the idempotent
defining that retract.  It follows that the unit gives an isomorphism
of the integral bottom-row computing complexes
\[
 \Hom_\Lambda(P_\bullet,\underline{\Z_p})
       \xrightarrow{\ \sim\ }
 \Hom_\Lambda(P_\bullet,H^0(K_n)).
\]
After $[1/p]$, \eqref{eq:group-cochain-bridge} and
\eqref{eq:geometric-group-descent} identify the restricted group
class with the constants bottom-row edge class induced by
$u_{\Omega_C^{n-1},U,\mathrm b}$.  This uses only the point-value
isomorphism, not an identification of the entire condensed module.

For a fixed $j$, the rationalized row is
computed by
\[
 \Hom_\Lambda(P_\bullet,H^j(K_n))[1/p].
\]
Every term is a direct summand of a finite sum of
$H^j(K_n)(*)[1/p]$.  By the integral Drinfeld theorem,
\cite[Theorem 1.1]{CDN2021}, a Galois element $\gamma$ acts on the latter
by the scalar $\chi_{\mathrm{cyc}}(\gamma)^{-j}$.
It commutes with $\Lambda$ and with the idempotents defining the
projective summands, so it has the same scalar action on the entire
row and on its cohomology.  No identification of a solid module by
its point value is involved.

The constants row $j=0$ has trivial Galois action.
Choose $\gamma\in G_F$ with $\chi_{\mathrm{cyc}}(\gamma)$ of infinite
order.  The only possible incoming differentials to bidegree $(3,0)$
are
\[
 d_2:E_2^{1,1}\longrightarrow E_2^{3,0},
 \qquad
 d_3:E_3^{0,2}\longrightarrow E_3^{3,0}.
\]
The second source is absent when $n=2$.
After inversion of $p$, the respective source scalars are
$\chi_{\mathrm{cyc}}(\gamma)^{-1}$ and
$\chi_{\mathrm{cyc}}(\gamma)^{-2}$, whereas the target scalar is $1$.
Every differential commutes with $\gamma$, so both are zero.
There are no outgoing differentials from the bottom row.
Hence the nonzero bottom-row image of the restriction of $e_{J_n}$
survives in $H^3_{\mathrm b}([\Omega_C^{n-1}/U])$.
Its abutment image is
$u_{\Omega_C^{n-1},U,\mathrm b}(\operatorname{res}^{J_n}_U e_{J_n})$
by the unit calculation
above.  Naturality \eqref{eq:constant-cochain-naturality} for
$U\hookrightarrow J_n$ identifies it with the pullback of
$u_{\Omega_C^{n-1},J_n,\mathrm b}(e_{J_n})$.
Thus the class on $[\Omega_C^{n-1}/J_n]=\mathcal M_n$ is nonzero.

Finally, that class is $G_F$-invariant.  Its classifying map to
$B_CJ_n$ is equivariant for the action on the $\Spd C$ base, while
$G_F$ fixes $J_n$ and the untwisted finite coefficients.  Naturality
of $u_{B_C,J_n,m}$ for this base action passes through the derived
coefficient limit and rationalization.  It therefore sends the
invariant group class to an invariant geometric class.
Equation~\eqref{eq:invariant-line} therefore places both nonzero
classes on the same one-dimensional $\Q_p$-line, proving
\cref{eq:primitive-comparison}.
\end{proof}

\subsection{Transport along the rank-one filtration}

\begin{theorem}\label{thm:geometric-transport}
There is a scalar $a\in\Q_p^\times$ such that the classes from the two
stabilizer factors satisfy
\begin{equation}\label{eq:geometric-transport}
 e_{P_3}|_{Y_C}=a\,e_{P_2}|_{Y_C}
       \quad\text{in }H^3_{\mathrm b}(Y_C).
\end{equation}
\end{theorem}

\begin{proof}
By \cref{lem:parabolic}, for the inclusion $\rho_3:D\to J_3$ and
projection $\rho_2:D\to J_2$ there is a scalar $b\in\Q_p^\times$
with
\[
 \rho_3^*e_{J_3}=b\,\rho_2^*e_{J_2}
       \quad\text{in }H^3_{\mathrm b}(D).
\]
Apply $u_{B_C,D,\mathrm b}$ to this group-cohomology equality and
pull it back along $d:\mathcal T\to B_CD$.
The exact filtration in \cref{prop:filtered-modification} identifies
the two $J_n$-classifying maps with the composites induced by
$D\to J_n$, and \eqref{eq:constant-cochain-naturality} identifies
their pullbacks with these same group classes.  The $P_n$-framing
torsors in that proposition give the first and last equalities below.
Consequently
\[
 \begin{aligned}
 \pi^*(e_{P_3}|_{Y_C})
   &=m_3^*(e_{P_3}|_{\mathcal M_3})\\
   &=\kappa_3\,d^*\rho_3^*e_{J_3}\\
   &=\kappa_3b\,d^*\rho_2^*e_{J_2}\\
   &=\frac{\kappa_3b}{\kappa_2}\,
                  m_2^*(e_{P_2}|_{\mathcal M_2})\\
   &=\frac{\kappa_3b}{\kappa_2}\,\pi^*(e_{P_2}|_{Y_C}).
 \end{aligned}
\]
All three scalars are nonzero.  The projective morphism
$\pi:\mathcal T\to Y_C$ satisfies
\cref{lem:projective-bundle}, so $\pi^*$ is injective.
Taking $a=\kappa_3b/\kappa_2$ proves the theorem.
\end{proof}

This transports the primitive through a filtration of rational local
systems.  The equality itself does not assert nonvanishing on $Y_C$.
The next section tests its coefficient image using the completed
tower and proves the required nonvanishing there.
